%% ==========================================================================
%%  SECTION 9 --- OPEN PROBLEMS
%%
%%  Reframed from the superseded sections/gap_and_direction.tex, which was
%%  written as a pitch for our own project. This is a SURVEY section: it
%%  reports the field's open problems, of which our interest is one item,
%%  named briefly at the end and not dwelt on.
%%
%%  EVIDENCE: an "open problem" is only listed here if a surveyed paper says
%%  it is open, or if the absence is visible in Table~\ref{tab:taxonomy}.
%%  Do not invent open problems.
%% ==========================================================================
\section{Open Problems}
\label{sec:open}

\todo{One framing paragraph: these are drawn from limitations the surveyed
      authors state themselves, and from gaps visible in
      Table~\ref{tab:taxonomy}. Say so --- it is what separates a survey's
      open-problems section from speculation.}

\subsection{Malicious security at scale}
\label{sec:open:malicious}

\todo{Nearly every system evaluated at realistic scale is semi-honest.
      nudge2026 lists malicious security as its first open question;
      grotto2023 notes that pika2022 generalised its predecessor technique to
      the malicious setting at a cost. Report what the papers say the barrier
      is, without asserting a solution.}

\subsection{Reducing the non-collusion assumption}
\label{sec:open:collusion}

\todo{The recurring trade: the fastest protocols assume two or three
      non-colluding servers. nudge2026 explicitly poses serving private
      recommendations with only two non-colluding parties as open. Contrast
      with the single-server line (simplepir2023, spiral2022), which removes
      the assumption at a computational cost. Present as a trade-off the field
      has not resolved, not as a defect.}

\subsection{The delivery stage}
\label{sec:open:delivery}

\todo{The sparse column in Table~\ref{tab:taxonomy}. nudge2026 states in its
      non-goals that it relies on ``other means'' --- private relays, Tor, or
      cryptographic PIR --- for users to fetch items privately. pirsona2021 is
      one of the few systems to address collection, training and delivery
      together. Meanwhile a substantial private-retrieval literature has
      developed largely separately (section 6).
      Report this as an observation about the literature; it is the gap our
      own project engages with, but this section is not the place to argue
      for that.}

\subsection{Leakage through the recommendation itself}
\label{sec:open:functionality}

\todo{The deepest open problem, and worth stating carefully because it is not
      a protocol failure. calandrino2011 attacks the recommender's INTENDED
      OUTPUT: no protocol flaw is involved. A system can be
      cryptographically perfect and still leak, because the functionality
      itself is revealing. asharov2018 make a structurally similar point for
      genomic similarity search.
      Cryptography bounds who sees the computation; it does not bound what the
      output implies. DP addresses this (mcsherry2009, and the DP extension in
      nudge2026) at a utility cost. Note that composing cryptographic and
      differential privacy in recommendation remains only partly explored.}

\subsection{Evaluation and comparability}
\label{sec:open:evaluation}

\todo{Forward from section 8: the field lacks a shared benchmark. Papers use
      different data sets, hardware, network models and quality metrics, so
      published numbers are frequently not comparable. This is a tractable,
      unglamorous open problem and worth naming as one.}

\subsection{Scope of our own follow-on work}
\label{sec:open:ours}

\todo{THREE OR FOUR SENTENCES, NO MORE. The course project that follows this
      survey engages with Section~\ref{sec:open:delivery}. Do not pitch, do
      not claim novelty, do not promise results --- state the interest and
      stop. Overselling here would undercut the survey's credibility.}
